\documentclass[12pt]{article}
\usepackage[T1]{fontenc}
\usepackage[utf8]{inputenc}
\usepackage[polish]{babel}
\usepackage{geometry}
\geometry{margin=1in}
\usepackage{listings}
\usepackage{xcolor}
\usepackage{graphicx}
\usepackage{hyperref}
\usepackage{amsmath}
\usepackage{amssymb}
\usepackage{tcolorbox}
\usepackage{framed}

\lstset{
    language=Python,
    basicstyle=\ttfamily\small,
    keywordstyle=\color{blue},
    commentstyle=\color{green!80!black},
    stringstyle=\color{red},
    backgroundcolor=\color{yellow!8},
    breaklines=true,
    showstringspaces=false,
    tabsize=4,
    captionpos=b,
    numbers=left,
    numberstyle=\tiny\color{black},
    stepnumber=1,
    numbersep=5pt,
    frame=single,
    xleftmargin=10pt,
    xrightmargin=10pt
}

\title{Predykcja złośliwości guzów nerek\\przy użyciu cech radiomicznych\\i algorytmów uczenia maszynowego}
\author{Piotr Przypaśniak}
\date{2025}

\begin{document}

\begin{framed}
\begin{center}
    {\Large Predykcja złośliwości guzów nerek\\
    przy użyciu cech radiomicznych i algorytmów uczenia maszynowego}
\end{center}
\end{framed}

\maketitle

\section{Wstęp}

Radiomika jest dziedziną medycyny zajmującą się ekstrakcją ilościowych cech z medycznych obrazów diagnostycznych (CT, MRI). Cechy radiomiczne zawierają informacje o teksturze, kształcie i intensywności obrazu, które mogą być predyktorami stanów patologicznych. W przypadku guzów nerek, predykcja złośliwości jest kluczowa dla planowania leczenia pacjenta.\\\\

Niniejszy projekt implementuje kompleksowy system klasyfikacji binarnej, którego celem jest predykcja tego, czy guz nerek jest złośliwy (malignant) czy łagodny (benign). System wykorzystuje trzy konkurencyjne modele uczenia maszynowego: Support Vector Machine (SVM), Random Forest oraz XGBoost.

\section{Architektura systemu}

System opiera się na klasycznym pipeline'u uczenia maszynowego, który można podzielić na cztery główne etapy: przetwarzanie danych, uczenie modelu, ewaluację oraz wydajność predykcyjną. Poniżej opisano każdy z komponentów.

\subsection{Przetwarzanie danych}

Moduł \texttt{preprocess.py} odpowiada za załadowanie i przygotowanie danych surowych. Główna funkcja to \textit{build\_original\_dataset}, która wykonuje następujące operacje:

\begin{enumerate}
    \item \textbf{Załadowanie danych} – wczytanie pliku CSV z cechami radiomicznymi ekstrahowanymi z obrazów medycznych.
    \item \textbf{Weryfikacja kolumny docelowej} – sprawdzenie obecności kolumny \textit{malignant}, która stanowi etykietę do klasyfikacji.
    \item \textbf{Selekcja cech} – filtrowanie kolumn, których nazwy zaczynają się od prefiksu \texttt{original\_}. Te kolumny reprezentują cechy radiomiczne.
    \item \textbf{Obsługa brakujących wartości} – usunięcie wierszy, w których brakuje etykiety docelowej, oraz obsługa niekompletnych rekordów (zachowywane są wiersze z co najmniej 90\% kompletnych danych).
    \item \textbf{Zapis przetworzonych danych} – zapisanie przygotowanego zbioru do pliku CSV.
\end{enumerate}

\noindent
Funkcja zwraca ramkę danych (DataFrame) zawierającą metadane (identyfikator przypadku, id annotator, numer guza, indeks wycinka) oraz wyekstrahowane cechy radiomiczne.

\begin{tcolorbox}[colframe=brown!75!black, colback=brown!10, title={\large Funkcja budowania zestawu danych}, fonttitle=\bfseries]
Moduł \texttt{src/preprocess.py} (linie 22–130) definiuje funkcję \textit{build\_original\_dataset}, której celem jest transformacja surowych danych radiomicznych w gotowy zbiór uczący. Kluczowe operacje to filtrowanie cech o prefiksie \texttt{original\_} oraz usunięcie wierszy z brakującymi wartościami.
\end{tcolorbox}

\begin{lstlisting}[caption={\texttt{src/preprocess.py} (linie 22--88): Budowanie i czyszczenie zestawu danych}]
# Linia 22-24: Załadowanie surowych danych
df = load_raw()

# Linia 30-31: Weryfikacja kolumny docelowej
if "malignant" not in df.columns:
    raise ValueError("Column 'malignant' not found")

# Linia 64-67: Selekcja cech radiomicznych
feature_columns = [
    c for c in df.columns
    if c.startswith("original_")
]

print(f"Found {len(feature_columns)} radiomics features")

# Linia 74-75: Obsługa braków danych - etykiety
df = df.dropna(subset=["malignant"])

# Linia 82-87: Obsługa braków danych - cechy
min_non_null = int(len(feature_columns) * 0.9)
df = df.dropna(
    thresh=min_non_null,
    subset=feature_columns
)
\end{lstlisting}

\subsection{Konfiguracja parametrów modeli}

Plik \texttt{config.py} przechowuje hiperparametry dla wszystkich trzech modeli. Wybrane parametry zostały dobrane empirycznie:

\begin{lstlisting}[caption={\texttt{src/config.py} (linie 9--29): Hiperparametry modeli klasyfikacyjnych}]
MODEL_PARAMS = {
    "SVM": {
        "C": 3,
        "kernel": "rbf",
        "gamma": "scale",
        "probability": True
    },
    "RandomForest": {
        "n_estimators": 500,
        "max_depth": 12,
        "class_weight": "balanced"
    },
    "XGBoost": {
        "n_estimators": 500,
        "max_depth": 5,
        "learning_rate": 0.02,
        "subsample": 0.9,
        "colsample_bytree": 0.9,
        "eval_metric": "logloss"
    }
}
\end{lstlisting}

\noindent
\begin{tcolorbox}[colframe=blue!75!black, colback=blue!10, title={\large Wybór hiperparametrów}, fonttitle=\bfseries]
\textit{SVM} z jądrem RBF (Radial Basis Function) jest przydatny dla wysokowymiarowych danych radiomicznych. Parametr $C=3$ kontroluje dywersyfikację błędów. \textit{RandomForest} z 500 drzewami zapewnia robustność poprzez agregację, zaś ograniczenie głębokości do 12 zapobiega overfittingowi. \textit{XGBoost} łączy gradient boosting z regularyzacją, gdzie mały learning rate (0.02) i subsample (0.9) zapewniają stabilność na złożonych zbiorach danych.
\end{tcolorbox}

\subsection{Modele klasyfikacyjne}

Moduł \texttt{models.py} dostarcza fabrykę do tworzenia instancji modeli na podstawie nazwy. Funkcja \textit{get\_model} enkapsuluje logikę inicjalizacji:

\begin{lstlisting}[caption={\texttt{src/models.py} (linie 8--18): Fabryka modeli}]
def get_model(name: str):
    params = MODEL_PARAMS[name]

    if name == "SVM":
        return SVC(**params)
    elif name == "RandomForest":
        return RandomForestClassifier(**params)
    elif name == "XGBoost":
        return XGBClassifier(**params)
    else:
        raise ValueError(f"Unknown model: {name}")
\end{lstlisting}

\noindent
Ta architektura ułatwia dodanie nowych modeli bez modyfikacji kodu treningowego. Każda instancja modelu jest inicjalizowana z hiperparametrami zdefiniowanymi w \texttt{config.py}.

\section{Pipeline uczenia maszynowego}

\subsection{Przygotowanie danych}

Etap uczenia w module \texttt{train.py} implementuje zaawansowany pipeline z wieloma etapami przetwarzania:

\begin{enumerate}
    \item \textbf{Imputation} (SimpleImputer) – uzupełnianie brakujących wartości medianą
    \item \textbf{Normalizacja} (StandardScaler) – standaryzacja cech do średniej 0 i wariancji 1
    \item \textbf{Filtr korelacji} (CorrelationFilter) – usunięcie cech o wysokiej korelacji ($> 0.90$)
    \item \textbf{Klasyfikacja} – aplikacja wybranego modelu
\end{enumerate}

\noindent
Pipeline'u towarzyszy następujący kod z pliku \texttt{src/train.py} (linie 132--137):

\begin{lstlisting}[caption={\texttt{src/train.py} (linie 132--137): Definicja pipeline'u przetwarzania}]
pipe = Pipeline([
    ("imputer", SimpleImputer(strategy="median")),
    ("scaler", StandardScaler()),
    ("correlation", CorrelationFilter(0.90)),
    ("model", model),
])
\end{lstlisting}

\subsection{Filtr korelacji}

\begin{tcolorbox}[colframe=green!75!black, colback=green!10, title={\large Redukcja wymiarowości radiomicznej}, fonttitle=\bfseries]
Filtr korelacji jest własnym transformatorem (ang. custom transformer) implementującym interfejs scikit-learn. W radomicznych cechach często występują wysokie korelacje między zmiennymi, które mogą prowadzić do overfittingu. Transformator eliminuje cechy skorelowane powyżej progu 0.90, zachowując reprezentatywny podzbiór cech.
\end{tcolorbox}

\begin{lstlisting}[caption={\texttt{src/train.py} (linie 27--60): Klasa CorrelationFilter dla redukcji wymiarowości}]
class CorrelationFilter(BaseEstimator, TransformerMixin):

    def __init__(self, threshold=0.90):
        self.threshold = threshold

    def fit(self, X, y=None):
        X_df = pd.DataFrame(X)
        corr = X_df.corr().abs()

        # Znalezienie kolumn do usunięcia
        upper = corr.where(
            np.triu(np.ones(corr.shape), k=1).astype(bool)
        )

        to_drop = [
            col for col in upper.columns
            if any(upper[col] > self.threshold)
        ]

        self.keep_columns_ = [
            c for c in X_df.columns
            if c not in to_drop
        ]

        print(f"Correlation filter removed: {len(to_drop)} features")

        return self

    def transform(self, X):
        X_df = pd.DataFrame(X)
        return X_df[self.keep_columns_]
\end{lstlisting}

\subsection{Walidacja krzyżowa z grupami}

Kluczowym aspektem jest \textit{StratifiedGroupKFold}, która zapewnia, że przypadki pacjentów nie są podzielone między zbiory treningowe i walidacyjne. Jest to ważne, ponieważ wiele wycieków danych z tego samego pacjenta może prowadzić do nierealistycznie wysokich wyników.

\begin{lstlisting}[caption={\texttt{src/train.py} (linie 114--147): Walidacja krzyżowa z group-aware splitting}]
# Linie 114-118: Definicja schematu CV
cv = StratifiedGroupKFold(
    n_splits=5,
    shuffle=True,
    random_state=RANDOM_STATE,
)

# Linie 139-147: Cross-validation z numeracją grup
scores = cross_val_score(
    pipe,
    X_train,
    y_train,
    cv=cv,
    groups=groups.loc[train_mask],
    scoring="roc_auc",
    n_jobs=-1
)
\end{lstlisting}

\noindent
\begin{tcolorbox}[colframe=purple!75!black, colback=purple!10, title={\large Zapobieganie wyciekom danych}, fonttitle=\bfseries]
\textit{StratifiedGroupKFold} jest niezbędna w kontekście medycznych danych radiomicznych. Bez separacji pacjentów między splitami, wiele wycieków tego samego przypadku tworzyłoby sztucznie zawyżone wyniki. Parametr \textit{groups} zawiera \texttt{case\_id}, zapewniając, że wszystkie próbki jednego pacjenta trafiają do tej samej części walidacyjnej.
\end{tcolorbox}

\subsection{Selekcja modelu}

System iteruje przez wszystkie dostępne modele, trenuje każdy przy użyciu walidacji krzyżowej i wybiera ten z najwyższym średnim wynikiem ROC AUC. Pełna logika znajduje się w \texttt{src/train.py} (linie 128--165):

\begin{lstlisting}[caption={\texttt{src/train.py} (linie 124--165): Pętla selekcji modelu}]
# Linie 124-126: Inicjalizacja zmiennych śledzących
best_score = -np.inf
best_name = None
best_pipe = None

# Linie 128-157: Iteracja po modelach
for name in MODEL_PARAMS.keys():
    model = get_model(name)

    pipe = Pipeline([
        ("imputer", SimpleImputer(strategy="median")),
        ("scaler", StandardScaler()),
        ("correlation", CorrelationFilter(0.90)),
        ("model", model),
    ])

    scores = cross_val_score(
        pipe, X_train, y_train,
        cv=cv,
        groups=groups.loc[train_mask],
        scoring="roc_auc",
        n_jobs=-1
    )

    mean_score = scores.mean()
    std_score = scores.std()

    print(f"{name}: ROC AUC = {mean_score:.3f} ± {std_score:.3f}")

    # Linie 154-157: Aktualizacja najlepszego modelu
    if mean_score > best_score:
        best_score = mean_score
        best_name = name
        best_pipe = pipe

# Linie 161-166: Finalny trening na całym zbiorze treningowym
best_pipe.fit(X_train, y_train)

print(f"\nBest model: {best_name}")
print(f"CV ROC AUC: {best_score:.3f}")
\end{lstlisting}

\section{Ewaluacja modelu}

Po wyborze najlepszego modelu, moduł \texttt{evaluate.py} przeprowadza kompleksową ewaluację na zbiorze testowym:

\begin{lstlisting}[caption={\texttt{src/evaluate.py} (linie 11--39): Funkcja ewaluacji}]
def evaluate(pipe, X_test, y_test, model_name: str):
    # Linie 12-16: Predykcja i metryki podstawowe
    y_pred = pipe.predict(X_test)

    print(f"\n=== Evaluation for {model_name} ===")
    acc = accuracy_score(y_test, y_pred)
    print(f"Accuracy: {acc:.3f}")

    # Linie 18-33: Metryki i wizualizacja dla modeli z predict_proba
    model = pipe.named_steps["model"]
    if hasattr(model, "predict_proba"):
        y_proba = pipe.predict_proba(X_test)[:, 1]
        auc = roc_auc_score(y_test, y_proba)
        print(f"ROC AUC: {auc:.3f}")

        fpr, tpr, _ = roc_curve(y_test, y_proba)
        plt.figure()
        plt.plot(fpr, tpr, label=f"AUC = {auc:.3f}")
        plt.plot([0, 1], [0, 1], linestyle="--")
        plt.xlabel("False Positive Rate")
        plt.ylabel("True Positive Rate")
        plt.title(f"ROC curve – {model_name}")
        plt.legend()
        plt.tight_layout()
        plt.show()

    # Linie 35-39: Raporty szczegółowe
    print("\nConfusion matrix:")
    print(confusion_matrix(y_test, y_pred))

    print("\nClassification report:")
    print(classification_report(y_test, y_pred))
\end{lstlisting}

\noindent
\begin{tcolorbox}[colframe=red!75!black, colback=red!10, title={\large Metryki ewaluacji}, fonttitle=\bfseries]
Ewaluacja obejmuje:
\begin{itemize}
    \item \textbf{Dokładność} (Accuracy) – odsetek poprawnych predykcji
    \item \textbf{ROC AUC} – pole pod krzywą ROC, odporne na niezrównoważone klasy
    \item \textbf{Macierz błędów} – liczba TP, TN, FP, FN
    \item \textbf{Raport klasyfikacji} – precyzja, recall, F1-score dla każdej klasy
    \item \textbf{Krzywa ROC} – wizualizacja wydajności klasyfikatora
\end{itemize}
\end{tcolorbox}

\section{Pipeline główny}

Moduł \texttt{main.py} orchestruje cały przepływ:

\begin{lstlisting}[caption={\texttt{src/main.py} (linie 1--6): Punkt wejścia aplikacji}]
# Linie 1-2: Import modułów
from src.train import train_all
from src.evaluate import evaluate

# Linie 4-6: Główna pętla
if __name__ == "__main__":
    pipe, X_test, y_test, best_name = train_all()
    evaluate(pipe, X_test, y_test, best_name)
\end{lstlisting}

\noindent
Funkcja \texttt{train\_all()} zwraca:
\begin{itemize}
    \item \texttt{pipe} – wytrenowany pipeline z najlepszym modelem
    \item \texttt{X\_test, y\_test} – zbiór testowy
    \item \texttt{best\_name} – nazwa wybranego modelu
\end{itemize}

\section{Przepływ danych}

Poniżej przedstawiono diagram przepływu danych w systemie:

\begin{verbatim}
┌─────────────────────────────────────────────────────────────┐
│                   Surowe dane radiomiczne                    │
│                    (data/raw/original.csv)                   │
└────────────────────────┬────────────────────────────────────┘
                         │
                         ▼
        ┌────────────────────────────────────┐
        │    preprocess.build_original_dataset()    │
        │  - Wczytanie danych                │
        │  - Filtrowanie cech radiomicznych │
        │  - Obsługa braków danych           │
        │  - Zapis przetworzonych danych    │
        └────────────────────┬───────────────┘
                             │
                             ▼
            ┌─────────────────────────────────┐
            │   Przetworzone dane             │
            │  (data/processed/original_clean.csv)  │
            └────────────────┬────────────────┘
                             │
                             ▼
            ┌─────────────────────────────────┐
            │   train.train_all()             │
            │  - Podział na zbióry            │
            │  - Pipeline przetwarzania       │
            │  - Walidacja krzyżowa           │
            │  - Selekcja modelu              │
            │  - Trening finalny              │
            └────────────────┬────────────────┘
                             │
                             ▼
            ┌─────────────────────────────────┐
            │   Wytrenowany model             │
            │   + Zbiór testowy               │
            └────────────────┬────────────────┘
                             │
                             ▼
            ┌─────────────────────────────────┐
            │   evaluate.evaluate()           │
            │  - Predykcja na zbiorze test.   │
            │  - Obliczenie metryk            │
            │  - Generacja raportów           │
            │  - Wizualizacja wyników         │
            └─────────────────────────────────┘
\end{verbatim}

\section{Technologiczne komponenty}

\subsection{Biblioteki}

\begin{itemize}
    \item \textbf{scikit-learn} – implementacja algorytmów ML, preprocessing, pipeline, ewaluacja (wykorzystane w \texttt{train.py}, \texttt{evaluate.py})
    \item \textbf{XGBoost} – gradient boosting model (importowany w \texttt{models.py}, wiersz 3)
    \item \textbf{pandas} – przetwarzanie danych tabelarycznych (fundamentalne w \texttt{preprocess.py})
    \item \textbf{numpy} – operacje numeryczne (operacje na macierzach w \texttt{train.py})
    \item \textbf{matplotlib} – wizualizacja wyników (krzywe ROC w \texttt{evaluate.py})
\end{itemize}

\subsection{Struktura katalogów}

\begin{lstlisting}[language=bash, caption=Hierarchia plików projektu]
kidney-tumors-malignancy-predictor/
├── src/
│   ├── __init__.py               
│   ├── config.py                  
│   ├── preprocess.py                
│   ├── models.py                    # Fabryka modeli 
│   ├── train.py                     
│   ├── evaluate.py                  
│   └── main.py                      
├── data/
│   ├── raw/                         
│   └── processed/                  
├── examples/                       
│   ├── 01_pyradiomics_quickstart.py
│   ├── 02_pyradiomics_filters.py
│   ├── 03_pyradiomics_feature_maps.py
│   ├── 04_pandas_features_overview.py
│   ├── 05_sklearn_introduction.py
│   └── 06_shap_explainability.py
├── PROJEKT.tex                      # Niniejsza dokumentacja
└── envs/env-radiomics.yaml          # Zależności środowiska Conda
\end{lstlisting}

\section{Zaawansowane aspekty}

\subsection{Problem niezrównoważonych klas}

W danych medycznych często występuje niezrównoważenie klas (więcej przypadków łagodnych niż złośliwych). System adresuje to poprzez:

\begin{itemize}
    \item \textbf{Stratified sampling} – walidacja krzyżowa ze \textit{StratifiedGroupKFold} (linie 114--118 w \texttt{train.py})
    \item \textbf{class\_weight='balanced'} – w RandomForest i XGBoost (linie 19 w \texttt{config.py})
    \item \textbf{ROC AUC} – metryka odporna na niezrównoważenie klas (linia 145 w \texttt{train.py})
\end{itemize}

\subsection{Zapobieganie wyciekom danych}

\begin{tcolorbox}[colframe=orange!75!black, colback=orange!10, title={\large Krytyczne kwestie walidacji w radiomice}, fonttitle=\bfseries]
Wysokowymiarowe cechy radiomiczne z tego samego pacjenta są silnie skorelowane. Bez odpowiedniej ostrożności, wiele wycieków tego samego przypadku tworzy sztucznie zawyżone wyniki. System zapobiega temu poprzez:

\begin{itemize}
    \item \textbf{GroupKFold} (linie 95--109 w \texttt{train.py}) – podział wg. przypadków pacjentów na podstawie \texttt{case\_id}, nie indywidualnych próbek
    \item \textbf{Holdout split} (linie 104--108 w \texttt{train.py}) – testowy zbiór zawiera zupełnie inne przypadki niż zbiór treningowy
    \item \textbf{Fit na zbiorze treningowym} – transformatory (imputer, scaler, correlation filter) dopasowywane tylko do danych treningowych
\end{itemize}
\end{tcolorbox}

\section{Podsumowanie}

Projekt implementuje kompleksowy system predykcji złośliwości guzów nerek przy użyciu nowoczesnych technik uczenia maszynowego. Architektura systemu jest modułowa i rozszerzalna:

\begin{itemize}
    \item \textbf{Modułowość} – każdy komponent (preprocessing, modele, training, evaluation) jest niezależny i łatwo testowany
    \item \textbf{Rozszerzalność} – łatwe dodanie nowych modeli poprzez modyfikację \texttt{config.py} i \texttt{models.py}
    \item \textbf{Robustwo naukowe} – rigorous validation procedures (group-aware cross-validation, holdout testing) zapewniają, że wyniki są wiarygodne i transferowalne do rzeczywistych aplikacji medycznych
\end{itemize}

\begin{tcolorbox}[colframe=black!75!black, colback=gray!10, title={\large Kluczowe osiągnięcia}, fonttitle=\bfseries]
\begin{enumerate}
    \item \textbf{Automatyczna selekcja modelu} – system autonomicznie wybiera optymalny model (SVM, RandomForest lub XGBoost) na podstawie walidacji krzyżowej
    \item \textbf{Redukcja wymiarowości} – filtr korelacji eliminuje redundantne cechy radiomiczne, zmniejszając overfitting
    \item \textbf{Zapobieganie wyciekom danych} – group-aware splitting zapewnia realistyczne szacunki wydajności w kontekście medycznym
    \item \textbf{Kompleksowa ewaluacja} – raport zawiera dokładność, ROC AUC, macierz błędów, precision/recall i wizualizacje
\end{enumerate}
\end{tcolorbox}

\vspace{1cm}
\noindent
System jest gotowy do wdrożenia i może służyć jako punkt wyjścia dla bardziej zaawansowanych technik, takich jak feature engineering, ensemble methods czy deep learning.

\end{document}
